\subsection{HOMOGENEOUS PRINCIPAL SETS}

DEFINITION 7. Let G be a group. An operation of G on a set E is called simply transitive if there exists an element x of E such that the orbital mapping defined by x is a bijection. A set E, together with a simply transitive left action of G on E, is called a left homogeneous principal G-set (or left homogeneous principal set under G).

It amounts to the same to say that G operates freely and transitively on E, or also that there exists an element \(x \in E\) such that the orbital mapping defined by x is an isomorphism of the G-set G (where G operates by left translation) onto E; or also that the two following conditions are satisfied:

\begin{enumerate}
\def\labelenumi{(\roman{enumi})}
\item
  E is non-empty.
\item
  for all elements \(x\) and \(y\) of \(\mathbf{E}\) , there exists one and only one element \(\alpha \in \mathbf{G}\) such that \(\alpha x = y\) .
\end{enumerate}

Condition (ii) is also equivalent to the following condition:

\begin{enumerate}
\def\labelenumi{(\roman{enumi})}
\setcounter{enumi}{2}
\tightlist
\item
  the mapping \((\alpha, x) \mapsto (\alpha x, x)\) is a bijection of \(G \times E\) onto \(E \times E\) .
\end{enumerate}

We leave to the reader the task of defining right homogeneous principal G-sets.

Examples. (1) Let G operate on itself by left (resp. right) translation. Thus a left (resp. right) homogeneous principal G-set structure is defined on G, which is sometimes denoted by \(G_{s}\) (resp. \(G_{d}\) ).

\begin{enumerate}
\def\labelenumi{(\arabic{enumi})}
\setcounter{enumi}{1}
\item
  Let E be a homogeneous set under a commutative group G. If G operates faithfully on E, the latter is a homogeneous principal G-set.
\item
  Let E and F be two isomorphic sets with structures of the same species and let \(\operatorname{Isom}(\mathbf{E}, \mathbf{F})\) be the set of isomorphisms of E onto F (with the given structures). The group \(\operatorname{Aut}(\mathbf{E})\) of automorphisms of E (with the given structure) operates on \(\operatorname{Isom}(\mathbf{E}, \mathbf{F})\) on the right by the law \((\sigma, f) \mapsto f \circ \sigma\) and \(\operatorname{Isom}(\mathbf{E}, \mathbf{F})\) is a right homogeneous principal \(\operatorname{Aut}(\mathbf{E})\) -set. Similarly, the group \(\operatorname{Aut}(\mathbf{F})\) operates on \(\operatorname{Isom}(\mathbf{E}, \mathbf{F})\) on the left by the law \((\sigma, f) \mapsto \sigma \circ f\) and \(\operatorname{Isom}(\mathbf{E}, \mathbf{F})\) is a left homogeneous principal \(\operatorname{Aut}(\mathbf{F})\) -set.
\item
  *A homogeneous principal set under the additive group of a vector space is called an affine space (cf.~II, § 9, no. 1).*
\end{enumerate}

The group of automorphisms of the homogeneous principal G-set \(G_{s}\) (Example 1) is the group of right translations of G which is identified with \(G^{0}\) (no. 5, Proposition 5). Let E be a homogeneous principal G-set and a an element of E. The orbital mapping \(\omega_{a}\) defined by a is an isomorphism of the G-set \(G_{s}\) onto E. By transporting the structure an isomorphism \(\psi_{a}\) is derived of \(G^{0}\) onto Aut(E). It should be noted that \(\psi_{a}\) in general depends on a; more precisely, for \(\alpha \in G\) ,

\[
\psi_ {\alpha a} = \psi_ {a} \circ \operatorname{Int} _ {\mathbf {G} ^ {0}} (\alpha) = \psi_ {a} \circ \operatorname{Int} (\alpha^ {- 1}).\tag{3}
\]

For, writing \(\delta_{a}\) for the translation \(x\mapsto x\alpha\) on G,

\[
\omega_ {\alpha a} = \omega_ {a} \circ \delta_ {\alpha}
\]

and

\[
\psi_ {a} (x) = \omega_ {a} \circ \delta_ {x} \circ \omega_ {a} ^ {- 1}, \quad x \in G,
\]

whence

\[
\psi_ {\alpha a} (x) = \omega_ {a} \circ \delta_ {\alpha} \circ \delta_ {x} \circ \delta_ {\alpha} ^ {- 1} \circ \omega_ {a} ^ {- 1} = \omega_ {a} \circ \delta_ {\alpha^ {- 1} x \alpha} \circ \omega_ {a} ^ {- 1} = \psi_ {a} (\alpha^ {- 1} x \alpha).
\]

